% Options for packages loaded elsewhere
\PassOptionsToPackage{unicode}{hyperref}
\PassOptionsToPackage{hyphens}{url}
\documentclass[
  a4paper,11pt]{article}
\usepackage{xcolor}
\usepackage[margin=1in]{geometry}
\usepackage{amsmath,amssymb}
\usepackage{graphicx}
\setcounter{secnumdepth}{-\maxdimen} % remove section numbering
\usepackage{iftex}
\ifPDFTeX
  \usepackage[T1]{fontenc}
  \usepackage[utf8]{inputenc}
  \usepackage{textcomp}
\else
  \usepackage{unicode-math}
  \defaultfontfeatures{Scale=MatchLowercase}
  \defaultfontfeatures[\rmfamily]{Ligatures=TeX,Scale=1}
\fi
\usepackage{lmodern}
\ifPDFTeX\else
  \setmainfont[]{Times New Roman}
\fi
\IfFileExists{upquote.sty}{\usepackage{upquote}}{}
\IfFileExists{microtype.sty}{%
  \usepackage[]{microtype}
  \UseMicrotypeSet[protrusion]{basicmath}
}{}
\makeatletter
\@ifundefined{KOMAClassName}{%
  \IfFileExists{parskip.sty}{%
    \usepackage{parskip}
  }{%
    \setlength{\parindent}{0pt}
    \setlength{\parskip}{6pt plus 2pt minus 1pt}}
}{% if KOMA class
  \KOMAoptions{parskip=half}}
\makeatother
\usepackage{longtable,booktabs,array}
\newcounter{none}
\usepackage{calc}
\usepackage{etoolbox}
\makeatletter
\patchcmd\longtable{\par}{\if@noskipsec\mbox{}\fi\par}{}{}
\makeatother
\IfFileExists{footnotehyper.sty}{\usepackage{footnotehyper}}{\usepackage{footnote}}
\makesavenoteenv{longtable}
\usepackage{bookmark}
\IfFileExists{xurl.sty}{\usepackage{xurl}}{}
\urlstyle{same}
\usepackage{hyperref}
\hypersetup{
  hidelinks,
  pdfcreator={LaTeX via pandoc}}
\setlength{\emergencystretch}{3em}
\providecommand{\tightlist}{%
  \setlength{\itemsep}{0pt}\setlength{\parskip}{0pt}}
\usepackage{amsthm}
\newtheorem{theorem}{Theorem}[section]
\newtheorem{proposition}[theorem]{Proposition}
\newtheorem{lemma}[theorem]{Lemma}
\newtheorem{corollary}[theorem]{Corollary}
\theoremstyle{definition}
\newtheorem{definition}[theorem]{Definition}
\theoremstyle{remark}
\newtheorem{remark}[theorem]{Remark}
\newtheorem{observation}[theorem]{Observation}
\begin{document}
\section{Shape-Invariant Traveling Waves on a 4-Dimensional Integer
Lattice: Existence Conditions, Conservation Structure, and
Self-Consistent
Loop}\label{shape-invariant-traveling-waves-on-a-4-dimensional-integer-lattice-existence-conditions-conservation-structure-and-self-consistent-loop}

\textbf{Author}: Noriaki Kihara\\
\textbf{Affiliation}: WF System Co., Ltd.\\
\textbf{ORCID}: 0009-0004-6753-4020\\
\textbf{DOI}: https://doi.org/10.5281/zenodo.19731598\\
\textbf{Version}: v1.0\\
\textbf{Date}: April 2026

\begin{center}\rule{0.5\linewidth}{0.5pt}\end{center}

\subsection{Abstract}\label{abstract}

On an \(n\)-dimensional integer lattice \(\mathbb{Z}^n\), we
systematically characterize the conditions under which a wave that
strictly preserves its amplitude distribution \(|\psi|\) under
translation --- a shape-invariant traveling wave --- can consistently
exist, along three axes: spherical consistency, multiplicative closure,
and mode decomposability. The existence conditions (D1)--(D4) for
spherical standing waves are satisfied in perfect-square dimensions
\(n \in \{4, 9, 16, \ldots\}\), but when one simultaneously requires an
associative multiplicative structure (Frobenius' theorem) and explicit
mode counting (Jacobi's four-square formula), \(n = 4\) is identified as
the \textbf{unique dimension}. The shape-invariant traveling waves
defined on this lattice simultaneously satisfy individuality,
coexistence, crossing invariance, composition rules, and symmetry-group
classification as labeled objects. Extension to the tangent bundle
\(\mathbb{Z}^4 \times \mathbb{Z}^4\) embeds a discrete flow, and the
self-return loop via orthogonal interaction waves derives a speed bound
from Nyquist stability. Furthermore, we show that central projection
from a parent space is the unique geometric mechanism that introduces
nonlinear effects into a linear system without invoking the concept of
time, and we argue that these structures point to \textbf{de Sitter
spacetime \(\mathrm{dS}_4\)} as a geometric necessity.

\textbf{Keywords}: shape-invariant traveling wave, integer lattice,
Hurwitz quaternion order, Jacobi four-square formula, Nyquist stability,
central projection, de Sitter spacetime

\begin{center}\rule{0.5\linewidth}{0.5pt}\end{center}

\subsection{§1 Introduction and Scope of
Claims}\label{introduction-and-scope-of-claims}

This paper describes, as purely geometric, number-theoretic, and
algebraic facts, how the dimensional dependence of waves on an
\(n\)-dimensional integer lattice \(\mathbb{Z}^n\) that strictly
preserve their waveform under translation --- shape-invariant traveling
waves --- is determined.

The claims of this paper are limited to the following four points:

\begin{enumerate}
\def\labelenumi{\arabic{enumi}.}
\tightlist
\item
  The intersection of the existence conditions (D1)--(D4) for spherical
  standing waves and the structural conditions (F1)--(F5) for
  shape-invariant traveling waves identifies \(n = 4\) as the unique
  consistent dimension.
\item
  Shape-invariant traveling waves on \(\mathbb{Z}^4\) simultaneously
  satisfy the five conditions of individuality, coexistence, crossing
  invariance, composition rules, and symmetry-group classification as
  labeled objects.
\item
  Through tangent bundle extension and orthogonal interaction waves, a
  self-return loop is constructed, and the speed bound is geometrically
  derived from the Nyquist stability condition.
\item
  Central projection from a parent space is the unique geometric
  mechanism that introduces nonlinear effects into a linear system
  without invoking the concept of time.
\end{enumerate}

The following elements are \textbf{outside the scope of this paper}:

\begin{itemize}
\tightlist
\item
  Dynamical stability (eigenvalue analysis for time evolution, convexity
  of energy functionals)
\item
  Rigorous formulation as solutions of nonlinear wave equations
\item
  Identification with physical particles, reproduction of quantum field
  theory
\end{itemize}

The objects called ``shape-invariant traveling waves'' in this paper are
clearly distinguished from ``solitons,'' which carry dynamical
implications. ``Shape-invariant'' refers solely to the geometric
property that \(|\psi|\) is preserved under translation. ``Traveling''
refers only to covariance with respect to translation as a lattice
automorphism, and does not include dynamical concepts of motion or
propagation.

\begin{center}\rule{0.5\linewidth}{0.5pt}\end{center}

\subsection{§2 Number-Theoretic Privilege of the 4-Dimensional Integer
Lattice}\label{number-theoretic-privilege-of-the-4-dimensional-integer-lattice}

\subsubsection{§2.1 Lagrange's Four-Square Theorem and Completeness of
the Squared-Radius
Set}\label{lagranges-four-square-theorem-and-completeness-of-the-squared-radius-set}

Consider the set of squared distances between lattice points in
\(\mathbb{Z}^n\): \(\{|\mathbf{x}|^2 : \mathbf{x} \in \mathbb{Z}^n\}\).

For \(n = 2\), values such as \(3, 6, 7, \ldots\) are missing, and for
\(n = 3\), integers of the form \(4^k(8m+7)\) are missing. By Lagrange's
four-square theorem, only for \(n \geq 4\) does

\[\{|\mathbf{x}|^2 : \mathbf{x} \in \mathbb{Z}^n\} = \mathbb{Z}_{\geq 0}\]

hold. \(n = 4\) is the \textbf{minimum dimension} at which this
completeness is achieved.

\subsubsection{§2.2 Jacobi's Four-Square Formula and Non-Emptiness of
Spherical Lattice
Points}\label{jacobis-four-square-formula-and-non-emptiness-of-spherical-lattice-points}

For the number \(r_n(N)\) of \(\mathbb{Z}^n\) lattice points on the
hypersphere of radius \(\sqrt{N}\), Jacobi's four-square formula gives

\[r_4(N) = 8 \sum_{\substack{d \mid N \\ 4 \nmid d}} d > 0 \quad (N \geq 1)\]

\(r_4(N)\) is positive for all \(N \geq 1\) and, moreover, possesses a
\textbf{closed explicit formula} in terms of divisor-sum functions.

\(r_3(N)\) has zeros at \(N = 4^k(8m+7)\), and the explicit formulas for
\(r_5, r_6, r_7\) involve cuspidal parts of automorphic forms and are
far more complex. The property of being ``positive for all \(N\) and
expressible in terms of divisor sums'' holds only for \(r_4\).

\subsubsection{§2.3 Integrality of the Equal-Component
Direction}\label{integrality-of-the-equal-component-direction}

The length of the main diagonal of an \(n\)-dimensional unit hypercube
is \(\sqrt{n}\), and the length of a lattice segment along the
equal-component direction \((\pm k, \pm k, \ldots, \pm k)\) is
\(\sqrt{n} \cdot k\). \(\sqrt{n}\) is an integer only when \(n\) is a
perfect square, and \(n = 4\) (\(\sqrt{4} = 2\)) is the \textbf{smallest
non-trivial such dimension}.

\subsubsection{§2.4 Multiplicative Structure of the Hurwitz Quaternion
Order}\label{multiplicative-structure-of-the-hurwitz-quaternion-order}

Within the quaternion algebra \(\mathbb{H}\), taking

\[\mathcal{H} = \mathbb{Z}^4 \cup \left(\mathbb{Z}^4 + \left(\tfrac{1}{2}, \tfrac{1}{2}, \tfrac{1}{2}, \tfrac{1}{2}\right)\right)\]

yields a \textbf{left Euclidean domain} satisfying norm multiplicativity
\(|xy| = |x| \cdot |y|\).

By Frobenius' theorem, the only finite-dimensional associative division
algebras over the reals are \(\mathbb{R}, \mathbb{C}, \mathbb{H}\). The
multiplicative integer lattices are limited to
\(\mathbb{Z}, \mathbb{Z}[i], \mathcal{H}\), with dimensions \(1, 2, 4\).
\(n = 4\) is the \textbf{maximum dimension} admitting an associative
multiplicative structure.

\begin{center}\rule{0.5\linewidth}{0.5pt}\end{center}

\subsection{§3 Existence Conditions for Spherical Standing Waves
(D1)--(D4)}\label{existence-conditions-for-spherical-standing-waves-d1d4}

We decompose the conditions for standing waves consistent with the
family of spheres \(\{S(\sqrt{N})\}_{N \geq 0}\) on \(\mathbb{Z}^n\)
into the following four conditions.

\textbf{(D1) Integrality of nodal radii}: The radius sequence of the
spherical family forming nodes is arranged at integer intervals along
the equal-component direction. Requires \(\sqrt{n} \in \mathbb{Z}\).

\textbf{(D2) Non-emptiness of nodal spheres}: For any nodal radius,
lattice points exist on that sphere. Requires \(r_n(N) > 0\) for all
\(N \geq 1\).

\textbf{(D3) Completeness of the nodal-radius-squared set}: The range of
values taken by the squared nodal radii covers all of
\(\mathbb{Z}_{\geq 0}\). Requires
\(\{|\mathbf{x}|^2 : \mathbf{x} \in \mathbb{Z}^n\} = \mathbb{Z}_{\geq 0}\).

\textbf{(D4) Wavenumber duality}: The nodal lattice in real space and
wavenumber space share an isomorphic discrete structure (self-duality of
\(\mathbb{Z}^n\)).

\textbf{Dimensional dependence}:

{\def\LTcaptype{none} % do not increment counter
\begin{longtable}[]{@{}
  >{\raggedright\arraybackslash}p{(\linewidth - 10\tabcolsep) * \real{0.1833}}
  >{\raggedright\arraybackslash}p{(\linewidth - 10\tabcolsep) * \real{0.1167}}
  >{\raggedright\arraybackslash}p{(\linewidth - 10\tabcolsep) * \real{0.1167}}
  >{\raggedright\arraybackslash}p{(\linewidth - 10\tabcolsep) * \real{0.1167}}
  >{\raggedright\arraybackslash}p{(\linewidth - 10\tabcolsep) * \real{0.1167}}
  >{\raggedright\arraybackslash}p{(\linewidth - 10\tabcolsep) * \real{0.3500}}@{}}
\toprule\noalign{}
\begin{minipage}[b]{\linewidth}\raggedright
Condition
\end{minipage} & \begin{minipage}[b]{\linewidth}\raggedright
\(n=2\)
\end{minipage} & \begin{minipage}[b]{\linewidth}\raggedright
\(n=3\)
\end{minipage} & \begin{minipage}[b]{\linewidth}\raggedright
\(n=4\)
\end{minipage} & \begin{minipage}[b]{\linewidth}\raggedright
\(n=9\)
\end{minipage} & \begin{minipage}[b]{\linewidth}\raggedright
\(n \geq 5\) general
\end{minipage} \\
\midrule\noalign{}
\endhead
\bottomrule\noalign{}
\endlastfoot
(D1) & × (\(\sqrt{2}\)) & × (\(\sqrt{3}\)) & ✓ & ✓ & Perfect squares
only \\
(D2) & × & × & ✓ & ✓ & \(n \geq 4\) \\
(D3) & × & × & ✓ & ✓ & \(n \geq 4\) \\
(D4) & ✓ & ✓ & ✓ & ✓ & All \(n\) \\
\end{longtable}
}

\textbf{Proposition 3.1}: The set of dimensions simultaneously
satisfying (D1)--(D4) is
\(\{k^2 : k \in \mathbb{Z}_{\geq 2}\} = \{4, 9, 16, 25, \ldots\}\), and
the minimum element is \(n = 4\).

\begin{center}\rule{0.5\linewidth}{0.5pt}\end{center}

\subsection{§4 Definition and Structural Conditions for Shape-Invariant
Traveling Waves
(F1)--(F5)}\label{definition-and-structural-conditions-for-shape-invariant-traveling-waves-f1f5}

\subsubsection{§4.1 Definition of Shape-Invariant Traveling
Waves}\label{definition-of-shape-invariant-traveling-waves}

\textbf{Definition 4.1}: A function \(\psi : \Lambda \to \mathbb{C}\) on
a lattice \(\Lambda \subset \mathbb{R}^n\) is a \textbf{shape-invariant
traveling wave} if there exists a wavevector
\(\mathbf{k} \in \Lambda^*\) such that for any translation
\(\mathbf{a} \in \Lambda\),

\[\psi(\mathbf{x} + \mathbf{a}) = e^{i\langle \mathbf{k}, \mathbf{a} \rangle} \cdot \psi(\mathbf{x}), \quad \langle \mathbf{k}, \mathbf{a} \rangle \in \mathbb{Z}\]

holds.

This definition involves neither time, equations, nor energy, and is
described solely in terms of a lattice, an inner product, and an
exponential map.

\subsubsection{§4.2 Structural Conditions}\label{structural-conditions}

The geometric conditions for shape-invariant traveling waves to
``consistently exist'':

\textbf{(F1) Phase consistency}:
\(\langle \mathbf{k}, \mathbf{a} \rangle \in \mathbb{Z}\) for all
\(\mathbf{a} \in \Lambda\) \(\Leftrightarrow\)
\(\mathbf{k} \in \Lambda^*\)

\textbf{(F2) Amplitude preservation}:
\(|\psi(\mathbf{x} + \mathbf{a})| = |\psi(\mathbf{x})|\) (a direct
consequence of the definition)

\textbf{(F3) Composition closure}: \(\psi_1 \cdot \psi_2\) is again a
shape-invariant traveling wave \(\Leftrightarrow\)
\(\mathbf{k}_1 + \mathbf{k}_2 \in \Lambda^*\)

\textbf{(F4) Spherical consistency}: The magnitude of the wavevector
\(|\mathbf{k}|\) lies on the radius sequence of the spherical family
\(\Leftrightarrow\) (D1)--(D4)

\textbf{(F5) Multiplicative closure}:
\(\psi(\mathbf{x}) \cdot \psi(\mathbf{y})\) can be written as a wave on
the lattice product \(\mathbf{x} \cdot \mathbf{y}\) \(\Leftrightarrow\)
an associative multiplicative structure on \(\Lambda\)

\subsubsection{§4.3 Dimensional Sieve of the
Conditions}\label{dimensional-sieve-of-the-conditions}

\begin{itemize}
\tightlist
\item
  (F1)(F2)(F3): Hold for all \(n\) when \(\Lambda = \mathbb{Z}^n\),
  \(\Lambda^* = \mathbb{Z}^n\)
\item
  (F4): \(n \in \{4, 9, 16, 25, \ldots\}\) (perfect-square dimensions)
\item
  (F5): \(n \in \{1, 2, 4\}\) (\(\mathbb{R}, \mathbb{C}, \mathbb{H}\);
  the octonions \(n = 8\) are excluded due to non-associativity)
\end{itemize}

\textbf{Theorem 4.2}: \((F4) \cap (F5) = \{4\}\).

The integer lattice on which shape-invariant traveling waves
simultaneously satisfy spherical consistency and multiplicative closure
is \textbf{limited to} \(\mathbb{Z}^4\).

\begin{center}\rule{0.5\linewidth}{0.5pt}\end{center}

\subsection{§5 Conserved Quantities and Invariant Structures
(C1)--(C8)}\label{conserved-quantities-and-invariant-structures-c1c8}

From the definition of a shape-invariant traveling wave
\(\psi(\mathbf{x} + \mathbf{a}) = e^{i\langle \mathbf{k}, \mathbf{a} \rangle} \cdot \psi(\mathbf{x})\),
the following quantities are \textbf{strictly conserved} under
translation \(\mathbf{a} \in \mathbb{Z}^4\).

\textbf{(C1) Amplitude distribution}:
\(|\psi(\mathbf{x} + \mathbf{a})| = |\psi(\mathbf{x})|\)

\textbf{(C2) Nodal set configuration}:
\(\{\mathbf{x} : \psi(\mathbf{x}) = 0\}\) is an automorphism under
translation by \(\mathbf{a}\)

\textbf{(C3) Phase difference structure}:
\(\langle \mathbf{k}, \mathbf{x}_i - \mathbf{x}_j \rangle \in \mathbb{Z}\)
is preserved for any pair of lattice points

\textbf{(C4) Wavevector magnitude}: \(|\mathbf{k}|\) is invariant under
translation

\textbf{(C5) Mode multiplicity}: \(r_4(|\mathbf{k}|^2)\) is invariant
under translation

In four dimensions, additional invariant structures exist with respect
to symmetry groups beyond translation:

\textbf{(C6) Conservation under quaternion multiplication}: Under
\(\mathbf{x} \mapsto \mathbf{x} \cdot u\) by a unit element \(u\)
(\(|u| = 1\)) of the Hurwitz order, the spherical average of \(|\psi|\)
is preserved. This is described as an action of
\(\mathrm{SO}(4) \cong (\mathrm{SU}(2) \times \mathrm{SU}(2))/\mathbb{Z}_2\).

\textbf{(C7) Conservation under \(D_4\) triality}: The theta series
\(\theta_{D_4}(\tau)\) is invariant under cyclic permutation of the
vector, spinor, and conjugate spinor representations.

\textbf{(C8) Conservation under modular transformations}:
\(\theta_{\mathbb{Z}^4}(\tau)\) transforms as a modular form of weight 2
under \(\Gamma_0(4)\).

\textbf{Proposition 5.1}: (C1)--(C5) hold in all dimensions, but
(C6)--(C8) hold only for \(n = 4\). The shape-invariant traveling wave
in four dimensions is the \textbf{unique structure} that simultaneously
preserves invariants under translation, quaternion multiplication,
triple duality, and modular transformations.

\begin{center}\rule{0.5\linewidth}{0.5pt}\end{center}

\subsection{§6 Coexistence of Multiple Waves and Crossing Invariance
(G1)--(G5)}\label{coexistence-of-multiple-waves-and-crossing-invariance-g1g5}

\subsubsection{§6.1 Superposition and Complete
Separation}\label{superposition-and-complete-separation}

Consider multiple shape-invariant traveling waves
\(\psi_{\mathbf{k}_j}(\mathbf{x}) = e^{i\langle \mathbf{k}_j, \mathbf{x} \rangle}\)
on \(\mathbb{Z}^4\).

The operation of extracting each component from the superposition
\(\Psi(\mathbf{x}) = \sum_j c_j \cdot \psi_{\mathbf{k}_j}(\mathbf{x})\)
can be performed exactly by the uniqueness of the discrete Fourier
transform on \(\mathbb{Z}^4\):

\[c_j = \lim_{N \to \infty} \frac{1}{(2N+1)^4} \sum_{\mathbf{x} \in [-N,N]^4} \Psi(\mathbf{x}) \cdot e^{-i\langle \mathbf{k}_j, \mathbf{x} \rangle}\]

Since the dual basis of \(\mathbb{Z}^4\) forms a complete orthogonal
system, any finite or countable number of labeled waves are
\textbf{completely separated}.

\subsubsection{§6.2 Geometric Realization of Particle-Like
Conditions}\label{geometric-realization-of-particle-like-conditions}

Regarding a shape-invariant traveling wave
\(\psi_{\mathbf{k}}(\mathbf{x})\) as a ``geometric object identified by
label \(\mathbf{k}\),'' the following hold in four dimensions:

\textbf{(G1) Individuality}: Each \(\mathbf{k} \in \mathbb{Z}^4\)
provides a unique label, and objects are identified by a countable
discrete set.

\textbf{(G2) Coexistence}: Any number of labeled objects consistently
coexist in the same space (closure under superposition).

\textbf{(G3) Crossing invariance}: Even when multiple objects occupy the
same region, each label is \textbf{completely recovered} by Fourier
decomposition.

\textbf{(G4) Composition rule}: Composition is defined by label addition
\(\mathbf{k}_1 + \mathbf{k}_2 \in \mathbb{Z}^4\), and is closed as
multiplication in the Hurwitz order.

\textbf{(G5) Classification system}: Labels are organized
group-theoretically under \(\mathrm{SO}(4)\), \(D_4\) triality, and
modular transformations.

\textbf{Theorem 6.1}: The unique dimension in which shape-invariant
traveling waves can exist on an integer lattice while simultaneously
satisfying (G1)--(G5) is \(n = 4\).

\textbf{Remark}: (G1)--(G5) embed structural correspondences to the
geometric conditions required of particles, but this paper claims only
the fact that ``the geometric structure satisfies the conditions'' and
does not make the identification ``therefore they are particles.''

\begin{center}\rule{0.5\linewidth}{0.5pt}\end{center}

\subsection{§7 Label Structure (L1)--(L5)}\label{label-structure-l1l5}

The wavevector \(\mathbf{k} \in \mathbb{Z}^4\) of a shape-invariant
traveling wave functions as an integer-valued label. In four dimensions,
this label set simultaneously possesses the following five layers of
structure.

\textbf{(L1) Group structure}: The label set \(\mathbb{Z}^4\) is closed
as a quaternion additive group and further possesses the multiplicative
group structure of the Hurwitz order \(\mathcal{H}\). Addition and
multiplication between labels are group-theoretically closed.

\textbf{(L2) Geometric hierarchy}: \(|\mathbf{k}|^2\) takes integer
values, and the label space is naturally stratified into spherical
shells by the Jacobi four-square formula \(r_4(|\mathbf{k}|^2)\). The
number of labels on each shell \(S(\sqrt{N})\) is explicitly given by a
divisor sum.

\textbf{(L3) Symmetry group action}:
\(\mathrm{SO}(4) \cong (\mathrm{SU}(2) \times \mathrm{SU}(2))/\mathbb{Z}_2\)
acts on the label set, and labels are organized as representations of
the 4-dimensional rotation group.

\textbf{(L4) Triality}: \(D_4\) triality classifies labels into three
types of representations (vector, spinor, and conjugate spinor), which
are cyclically permutable.

\textbf{(L5) Modular transformations}: The theta series
\(\theta_{\mathbb{Z}^4}(\tau)\) transforms modularly under
\(\Gamma_0(4)\), establishing scale duality of the label distribution.

\textbf{Proposition 7.1}: The label structure simultaneously satisfying
all of (L1)--(L5) is limited to \(\mathbb{Z}^4\). (L1) holds for
\(n \in \{1, 2, 4\}\), (L2) admits a concise closed form only for
\(n = 4\), and (L4) holds only for \(n = 4\).

\begin{center}\rule{0.5\linewidth}{0.5pt}\end{center}

\subsection{§8 Tangent Bundle Extension and Discrete
Flow}\label{tangent-bundle-extension-and-discrete-flow}

\subsubsection{§8.1 Tangent Bundle Extension of the Label
Space}\label{tangent-bundle-extension-of-the-label-space}

We extend the label \(\mathbf{k} \in \mathbb{Z}^4\) of a shape-invariant
traveling wave to the pair of label and rate of change
\((\mathbf{k}, \Delta\mathbf{k}) \in \mathbb{Z}^4 \times \mathbb{Z}^4\).
This corresponds to the \textbf{discrete tangent bundle}
\(T\mathbb{Z}^4 = \mathbb{Z}^4 \times \mathbb{Z}^4\) of the label space
\(\mathbb{Z}^4\).

The extended wave is written as

\[\psi_{(\mathbf{k}, \Delta\mathbf{k})}(\mathbf{x}, t) = e^{i\langle \mathbf{k} + t \cdot \Delta\mathbf{k},\; \mathbf{x} \rangle}\]

Here \(t \in \mathbb{Z}\) is a parameter of the discrete flow and is not
physical time.

\subsubsection{§8.2 Discrete Flow}\label{discrete-flow}

The self-update map
\(U : (\mathbf{k}, \Delta\mathbf{k}) \mapsto (\mathbf{k} + \Delta\mathbf{k}, \Delta\mathbf{k})\)
generates a discrete flow on \(\mathbb{Z}^4 \times \mathbb{Z}^4\). The
wave itself embeds the rule \(\Delta\mathbf{k}\) for how to advance its
label, and the external concept of time appears only as the flow
parameter \(t\).

\subsubsection{§8.3 Structures Specific to Four
Dimensions}\label{structures-specific-to-four-dimensions}

The extended label space \(\mathbb{Z}^4 \times \mathbb{Z}^4\) possesses
the double structure \(\mathcal{H} \times \mathcal{H}\) of the Hurwitz
order. The self-update flow is organized as an \(\mathrm{SO}(4)\)
action, and the quaternion product operations on
\((\mathbf{k}, \Delta\mathbf{k})\) remain closed within
\(\mathbb{Z}^4 \times \mathbb{Z}^4\).

\begin{center}\rule{0.5\linewidth}{0.5pt}\end{center}

\subsection{§9 Self-Return Loop and Nyquist
Stability}\label{self-return-loop-and-nyquist-stability}

\subsubsection{§9.1 Emission of Orthogonal Interaction
Waves}\label{emission-of-orthogonal-interaction-waves}

From the central phase of a wave
\(\psi_{(\mathbf{k}, \Delta\mathbf{k})}\), an interaction wave \(\chi\)
is emitted in directions orthogonal to the wavevector \(\mathbf{k}\) in
four dimensions. By the spherical wave regularity of (D1)--(D4),
\(\chi\) is a 4-dimensional spherical wave that spreads in a
shape-invariant manner on \(\mathbb{Z}^4\).

\subsubsection{§9.2 Self-Recrossing
Condition}\label{self-recrossing-condition}

For the trajectory of the wave
\(\tau \mapsto \mathbf{x}_{\mathrm{self}}(\tau)\) and the current
position \(\mathbf{x}_{\mathrm{self}}(t)\), the recrossing condition

\[|\mathbf{x}_{\mathrm{self}}(t) - \mathbf{x}_{\mathrm{self}}(\tau)| = c(t - \tau)\]

is satisfied for some \(\tau < t\), meaning that a \(\chi_\tau\) emitted
in the past acts multiplicatively on the current \(\psi\). If the wave
is stationary (\(\Delta\mathbf{k} = 0\)), \(\chi\) merely spreads
outward and does not return to itself. Self-recrossing occurs only when
the wave is in motion (\(\Delta\mathbf{k} \neq 0\)), causing a label
transformation
\((\mathbf{k}, \Delta\mathbf{k}) \to (\mathbf{k}', \Delta\mathbf{k}')\).

\subsubsection{§9.3 Nyquist Stability and Speed
Bound}\label{nyquist-stability-and-speed-bound}

We analyze this self-return loop from a control-theoretic perspective.
The phase structure of the open-loop transfer function:

\begin{itemize}
\tightlist
\item
  \textbf{Fundamental phase}: \(-\pi/2\) (emission in the orthogonal
  direction)
\item
  \textbf{Delay phase}: \(-\omega \Delta t\) (recrossing delay)
\item
  \textbf{Composite phase}: \(-\pi/2 - \omega \Delta t\)
\end{itemize}

The amplitude decay of 4-dimensional spherical waves is \(1/r^{3/2}\),
and the loop gain decays sufficiently fast with distance. The Nyquist
trajectory does not encircle the \((-1, 0)\) point, and the loop is
stable.

\textbf{Derivation of the speed bound}: From the effective delay
\(\Delta t \propto 1/(c - v_\parallel)\), as the wave's travel speed
\(v\) approaches the propagation speed \(c\), we have
\(\Delta t \to 0\), and the divergence boundary frequency
\(\omega_c = \pi/(2\Delta t) \to \infty\). At \(v = c\), all frequencies
reach the divergence boundary.

\textbf{Proposition 9.1}: As a condition for maintaining Nyquist
stability, \(v \leq c\) is geometrically derived. The speed bound and
displacement bound are unified as different approach paths to the same
stability boundary \(\angle G = -\pi, |G| = 1\).

\begin{center}\rule{0.5\linewidth}{0.5pt}\end{center}

\subsection{§10 Central Projection and the Geometric Origin of Nonlinear
Effects}\label{central-projection-and-the-geometric-origin-of-nonlinear-effects}

\subsubsection{§10.1 Limitation of Linear
Systems}\label{limitation-of-linear-systems}

Within a single subjective space \(\mathbb{Z}^4\), shape-invariant
traveling waves \textbf{coexist without interaction} due to the closure
property of linear superposition. Addition operations alone cannot
transform the label set (by the uniqueness of Fourier decomposition).

\subsubsection{§10.2 Four Methods of Multiplicative
Action}\label{four-methods-of-multiplicative-action}

Methods for introducing nonlinear effects into a linear system:

\textbf{(i) Duality in phase space}: Formulate the system as a
\(U(1)\)-valued function and read phase addition as amplitude
multiplication.

\textbf{(ii) Transformation of the evaluation-point lattice}: Apply
quaternion rotation \(\mathbf{x} \mapsto \mathbf{x} \cdot u\) to the
evaluation-point space, inducing a multiplicative action
\(\mathbf{k} \mapsto \mathbf{k} \cdot u^*\) on the label space. Specific
to \(n = 4\).

\textbf{(iii) Phase composition of two systems}: Phase composition
through same-point evaluation of two independent systems.

\textbf{(iv) Central projection from a parent space}: Project a linear
wave on the parent space \(S^5(R)\) onto the child space via central
projection \(\pi : S^5 \to \mathbb{Z}^4\). Because the projection map is
a fractional (nonlinear) map, the result appears on the child space as a
function with \textbf{nonlinear phase}.

\subsubsection{§10.3 Privilege of Central
Projection}\label{privilege-of-central-projection}

Method (iv) is qualitatively different from the other three. Methods
(i)--(iii) exploit the internal structure of the child space
\(\mathbb{Z}^4\) and produce changes within the range of the child
space's symmetries. Method (iv) is an action from \textbf{outside the
frame}:

\begin{itemize}
\tightlist
\item
  Operations in the parent space are mere addition (linear)
\item
  On the child space, they appear as nonlinear changes of unknown origin
\item
  The symmetries of the child space themselves can be deformed ``from
  the outside''
\end{itemize}

\textbf{Proposition 10.1}: The unique method that introduces nonlinear
effects into a linear system in a geometrically justified manner without
introducing dynamics is central projection from a higher-dimensional
parent space.

\begin{center}\rule{0.5\linewidth}{0.5pt}\end{center}

\subsection{§11 Convergence to de Sitter
Spacetime}\label{convergence-to-de-sitter-spacetime}

\subsubsection{§11.1 Soliton Stability on Closed
Spheres}\label{soliton-stability-on-closed-spheres}

As shown in the companion paper {[}1{]}, shape-invariant standing waves
can exist stably on a closed \(n\)-dimensional sphere \(S^n\)
\textbf{only in odd dimensions}. On even-dimensional spheres, the hairy
ball theorem (the nonexistence of a nowhere-vanishing vector field)
causes the standing wave structure to break down.

\subsubsection{§11.2 Intersection of Two
Arguments}\label{intersection-of-two-arguments}

{\def\LTcaptype{none} % do not increment counter
\begin{longtable}[]{@{}
  >{\raggedright\arraybackslash}p{(\linewidth - 6\tabcolsep) * \real{0.1961}}
  >{\raggedright\arraybackslash}p{(\linewidth - 6\tabcolsep) * \real{0.2549}}
  >{\raggedright\arraybackslash}p{(\linewidth - 6\tabcolsep) * \real{0.4118}}
  >{\raggedright\arraybackslash}p{(\linewidth - 6\tabcolsep) * \real{0.1373}}@{}}
\toprule\noalign{}
\begin{minipage}[b]{\linewidth}\raggedright
Argument
\end{minipage} & \begin{minipage}[b]{\linewidth}\raggedright
Target space
\end{minipage} & \begin{minipage}[b]{\linewidth}\raggedright
Privileged dimension
\end{minipage} & \begin{minipage}[b]{\linewidth}\raggedright
Basis
\end{minipage} \\
\midrule\noalign{}
\endhead
\bottomrule\noalign{}
\endlastfoot
This paper (discrete side) & Infinite flat lattice \(\mathbb{Z}^n\) &
\(n = 4\) & (F4)∩(F5)=\{4\} \\
Companion paper (continuous side) & Closed sphere \(S^n\) & Odd
dimensions & Huygens' principle + hairy ball theorem \\
\end{longtable}
}

Both arguments hold independently, but they impose constraints on the
construction of physical spacetime:

\begin{itemize}
\tightlist
\item
  \textbf{Spatial part}: Solitons are stable on a closed sphere →
  \textbf{odd dimension} → minimum is \(S^3\)
\item
  \textbf{Temporal part}: Shape-invariant waves propagate along an
  infinitely open direction → \textbf{hyperbolic 1-dimension}
\item
  \textbf{Overall}: (D1)--(D4) hold in a locally flat 4-dimensional
  setting → \(3 + 1\) dimensions
\end{itemize}

\subsubsection{§11.3 Geometric Necessity of de Sitter
Spacetime}\label{geometric-necessity-of-de-sitter-spacetime}

De Sitter spacetime \(\mathrm{dS}_4\) is defined as the hyperboloid in
5-dimensional Minkowski space \(\mathbb{M}^{1,4}\):

\[-X_0^2 + X_1^2 + X_2^2 + X_3^2 + X_4^2 = R^2\]

The spatial part is a closed \(S^3\) (odd-dimensional, soliton-stable),
and the temporal part is hyperbolic (infinitely open, the arena for
shape-invariant wave propagation).

\textbf{Proposition 11.1}: The only 4-dimensional spacetime structure
that simultaneously satisfies the existence conditions for
shape-invariant traveling waves (discrete side: \(n = 4\) is unique) and
the existence conditions for closed solitons (continuous side: odd
dimensions only) is de Sitter spacetime \(\mathrm{dS}_4\), described as
the product structure of the space \(S^3\) and hyperbolic time.

\subsubsection{§11.4 Roles of Parent and Child
Spaces}\label{roles-of-parent-and-child-spaces}

Within the framework of the central projection series:

\begin{itemize}
\tightlist
\item
  \textbf{Child space} \(S^4(R)\) (even-dimensional, interior
  \(\mathbb{Z}^4\)): The arena for shape-invariant traveling waves.
  Closed solitons are unstable on the surface.
\item
  \textbf{Parent space} \(S^5(R)\) (odd-dimensional): The source of
  closed solitons. Introduces nonlinear effects into the child space via
  central projection.
\end{itemize}

The structure in which the parent space is odd-dimensional and the
spatial part of the child space is odd-dimensional is
\textbf{necessarily} determined by the geometric difference between odd
and even dimensions.

\begin{center}\rule{0.5\linewidth}{0.5pt}\end{center}

\subsection{§12 Summary}\label{summary}

We organize the logical flow presented in this paper:

\begin{enumerate}
\def\labelenumi{\arabic{enumi}.}
\tightlist
\item
  The existence conditions (D1)--(D4) for spherical standing waves on
  \(\mathbb{Z}^n\) are satisfied in perfect-square dimensions
  \(\{4, 9, 16, \ldots\}\) (§3).
\item
  Adding the structural conditions (F1)--(F5) for shape-invariant
  traveling waves, \((F4) \cap (F5) = \{4\}\) identifies \(n = 4\) as
  the unique dimension (§4).
\item
  Shape-invariant traveling waves on \(\mathbb{Z}^4\) possess conserved
  quantities (C1)--(C8), geometrically realize particle-like conditions
  (G1)--(G5), and are organized by label structures (L1)--(L5) (§5--§7).
\item
  Through tangent bundle extension \(\mathbb{Z}^4 \times \mathbb{Z}^4\),
  a discrete flow is defined, and the self-return loop of orthogonal
  interaction waves derives the speed bound \(v \leq c\) from Nyquist
  stability (§8--§9).
\item
  Central projection from a parent space is the unique geometric
  mechanism that introduces nonlinear effects into a linear system
  without invoking the concept of time (§10).
\item
  Two independent arguments from the discrete side (\(n = 4\)) and the
  continuous side (odd-dimensional spheres) point to de Sitter spacetime
  \(\mathrm{dS}_4\) as a geometric necessity (§11).
\end{enumerate}

All claims in this paper are stated as purely geometric,
number-theoretic, and algebraic facts, and require no physical
interpretation. Physical correspondences (particles, interactions,
spacetime structure) remain open as \textbf{interpretations connectable}
to these geometric facts, but do not affect the validity of the
propositions themselves.

\begin{center}\rule{0.5\linewidth}{0.5pt}\end{center}

\subsection{References}\label{references}

{[}1{]} Noriaki Kihara, ``Shape-Invariant Standing Waves on Closed
Spheres: Dimensional Stability and the Geometric Necessity of 3+1
Spacetime,'' 2026.

{[}2{]} Noriaki Kihara, ``A Three-Layer Model of the Universe via
Central Projection,'' Zenodo, 2025. DOI: 10.5281/zenodo.15083729.

{[}3{]} Noriaki Kihara, ``Regularity of Inter-Sphere Projection
(Gnomonic Orthoprojection),'' Zenodo, 2025. DOI:
10.5281/zenodo.15292757.

{[}4{]} Noriaki Kihara, ``(R,Q) Mapping Definition Paper {[}M3{]},''
Zenodo, 2025. DOI: 10.5281/zenodo.19692853.
\end{document}
